% answers_short.tex —— 2013 年·陕西卷·文科 · 简短答案速查
% 由 tools/make_short_answers.py 自 answers.tex 的 \gksolhead/\gkanswer 提取，
% 并与 2013/2013_数学真题_陕西文_答案要点.md 逐题比对后生成，勿手改（改源 MD 后重跑）。

\begin{gknotice}
\textbf{使用说明}：本节只给最终答案便于快速核对，完整推导见解析版。答案为编写者独立推导，\textbf{非官方标准答案}。
\end{gknotice}

\gkpart{选择题}
\begin{center}\begin{tabular}{ccccccccccc}
\toprule
题号 & 1 & 2 & 3 & 4 & 5 & 6 & 7 & 8 & 9 & 10 \\
\midrule
答案 & B & C & B & C & D & C & A & B & A & D \\
\bottomrule
\end{tabular}\end{center}

\gkpart{填空题}
\begin{itemize}[leftmargin=2.6em]
  \item[\textbf{11.}] $\dfrac54$。$a=4$，$b=3$，$c=\sqrt{a^2+b^2}=5$，故 $e=\dfrac ca=\dfrac54$。
  \item[\textbf{12.}] $3\pi$。三视图还原为半径 $1$ 的半球，表面积 $2\pi r^2+\pi r^2=3\pi$。
  \item[\textbf{13.}] $(n+1)(n+2)\cdots(n+n)=2^n\times1\times3\times\cdots\times(2n-1)$。左边为 $n$ 个连续整数之积，右边为 $2^n$ 与前 $n$ 个奇数之积。
  \item[\textbf{14.}] $20$。设矩形另一边为 $y$，由相似比 $\dfrac{x}{40}+\dfrac{y}{40}=1$ 得 $y=40-x$，面积 $x(40-x)$ 在 $x=20$ 时最大。
  \item[\textbf{15.}] $\mathbf R$。由绝对值不等式 $|x-a|+|x-b|\geqslant|a-b|>2$，故对任意实数 $x$ 都成立。
  \item[\textbf{16.}] $\sqrt6$。由 $\triangle AED\backsim\triangle CEB$ 得 $AE\cdot EB=DE\cdot CE$，结合 $PE\parallel BC$ 与 $PD=2DA=2$ 化得 $PE^2=6$。
  \item[\textbf{17.}] $(1,0)$。消去 $t$ 得 $y^2=4x$，其焦点为 $(1,0)$。
\end{itemize}

\gkpart{解答题}
\begin{itemize}[leftmargin=2.6em]
  \item[\textbf{18.}] (1) $T=\pi$。$f(x)=\sin\left(2x-\dfrac\pi6\right)$，故最小正周期为 $\pi$；(2) 最大值为 $1$，最小值为 $-\dfrac12$。在 $\left[0,\dfrac\pi2\right]$ 上 $2x-\dfrac\pi6\in\left[-\dfrac\pi6,\dfrac{5\pi}6\right]$
  \item[\textbf{19.}] (1) $S_n=\dfrac{n(a_1+a_n)}2=na_1+\dfrac{n(n-1)}2d$。由倒序相加或通项求和可得；(2) $\{a_n\}$ 是等比数列，$a_n=q^{n-1}$。由 $a_n=S_n-S_{n-1}$（$n\geqslant2$）得 $a_n=q^{n-1}$，且 $a_1=1$ 也符合
  \item[\textbf{20.}] (1) 平面 $A_1BD\parallel$ 平面 $CD_1B_1$，见证明；(2) $1$。$AO=1$，$A_1O=\sqrt{AA_1^2-AO^2}=1$，三棱柱底面积 $S_{\triangle ABD}=1$，高为 $1$，故体积为 $1$
  \item[\textbf{21.}] (1) $A$ 组 $3$ 人，$B$ 组 $6$ 人，$C$ 组 $9$ 人，$D$ 组 $9$ 人，$E$ 组 $3$ 人。抽样比为 $\dfrac6{100}=\dfrac3{50}$；(2) $\dfrac29$。$A$ 组抽 $3$ 人、$B$ 组抽 $6$ 人，各有 $2$ 人支持，故概率为 $\dfrac23\times\dfrac26=\dfrac29$
  \item[\textbf{22.}] (1) $\dfrac{x^2}4+\dfrac{y^2}3=1$。由 $|x-4|=2\sqrt{(x-1)^2+y^2}$ 平方整理即得；(2) $k=\pm\dfrac32$。设 $A(x_1,y_1)$，则 $B(2x_1,2y_1-3)$，两点均在椭圆上，解得 $y_1=\dfrac32$，$x_1=\pm1$
  \item[\textbf{23.}] (1) $y=x-1$。反函数为 $y=\ln x$，在 $x=1$ 处导数为 $1$，切线过 $(1,0)$；(2) 两曲线有唯一公共点 $(0,1)$，见证明；(3) $f\left(\dfrac{a+b}{2}\right)<\dfrac{f(b)-f(a)}{b-a}$。令 $t=\dfrac{b-a}2>0$，则右式为 $\mathrm e^{\frac{a+b}2}\cdot\dfrac{\sinh t}{t}$，而 $\dfrac{\sinh t}{t}>1$
\end{itemize}
